\documentclass[10pt]{article}
\usepackage[letterpaper,margin=0.85in]{geometry}
\usepackage[T1]{fontenc}\usepackage{lmodern}
\usepackage{amsmath}
\usepackage{graphicx,booktabs,tabularx,caption,subcaption,float,xcolor,enumitem,titlesec,multicol}
\usepackage[hidelinks]{hyperref}
\graphicspath{{figs/}}
\captionsetup{font=small,labelfont=bf,skip=4pt}
\titlespacing*{\section}{0pt}{7pt}{3pt}
\setlist{nosep,leftmargin=1.4em}
\renewcommand{\arraystretch}{1.0}

\title{\vspace{-16pt}\large 12-830 Principles of Digital Twins, Fall 2026 \\[2pt]
\Large Deliverable 2, Part 2.1: Physical System Characterization \\[2pt]
\normalsize\mdseries AttentionTwin: a context-aware digital twin of the driver--vehicle--environment system}
\author{\normalsize Chen Yu Chang, Hanwen Guo, Zihao Hu}
\date{\vspace{-6pt}\normalsize November 2026}

\begin{document}\maketitle\vspace{-16pt}

\noindent\textbf{Scope.} The physical system is one electric vehicle in SAE Level~2 supervised driving, with its driver and everything within its roughly 100\,m perception range (Deliverable~1). Three coupled states describe it: the \emph{driver} (where the eyes are, whether the hands are on the wheel, how alert), the \emph{vehicle} (speed, inputs, automation mode) and the \emph{environment} (traffic, pedestrians, road, weather). Safety depends on the match between the attention the driver supplies and the attention the scene demands. The twin estimates both and acts on the gap. Driver and vehicle states come from the car's own sensors; environment state comes from exterior perception on the road and, for development and calibration, from CARLA, where every agent's position is known exactly. The CARLA figures here show what the twin measures and why context changes the meaning of a glance.

\section{System layout and spatial relationships}
Figure~\ref{fig:layout} shows the geometry. In the cabin, one infrared camera at the rear-view mirror observes the driver's face and eyes from about 0.6\,m, a torque sensor in the steering column senses hands on the wheel, and seat and belt switches confirm occupancy. The vehicle's own state is read from the CAN bus. Outside, eight cameras and the perception stack see other road users, lanes and signals out to about 100\,m, the scale at which a pedestrian 18\,m ahead or a car closing at 2\,s headway is distinguishable from an empty road. Figure~\ref{fig:driver} is the driver's view with the gaze zones a cabin camera classifies, Figure~\ref{fig:scene} the environment state around the car at one moment, and Figure~\ref{fig:coupling} how the three states couple into the decisions.

\begin{figure}[H]\centering\includegraphics[width=\linewidth]{layout.png}
\caption{Layout of the physical system. Left: the vehicle and the environment within its perception range, with the quantity measured for each road user. Right: the cabin and vehicle sensing points.}\label{fig:layout}\end{figure}

\begin{figure}[H]\centering\includegraphics[width=0.9\linewidth]{coupling.png}
\caption{Coupling: driver, vehicle and environment states become attention supply and scene demand; their gap is the attention risk score that drives the fast (seconds) and slow (days) loops.}\label{fig:coupling}\end{figure}

\section{Physical processes}
\begin{enumerate}[label=\textbf{P\arabic*}]
\item \textbf{Glance behaviour and alertness.} Attention is a sequence of glances between zones (road, mirrors, cluster, screen, lap). Off-road glance duration, its accumulation over a window, blink rate and eyelid closure (PERCLOS) are the measurable signatures; glances over 2\,s roughly double crash risk and drowsiness raises it four to six times (Deliverable~1, [6]).
\item \textbf{Vehicle dynamics and automation mode.} Speed, longitudinal and lateral acceleration, steering and braking inputs, and whether Autosteer or FSD is engaged. Hard braking above 0.3\,g and turning above 0.4\,g are the thresholds Tesla's Safety Score already uses; the mode decides whether the driver is supervising or driving.
\item \textbf{Scene demand.} Traffic density, gap and time-to-collision to the lead vehicle, distance to the nearest pedestrian, junctions, signal state, speed limit, weather and light together set how much attention the moment requires. Figure~\ref{fig:timeseries} shows these quantities over a two-minute CARLA drive with an illustrative demand index built from them; Figure~\ref{fig:moments} shows the driver's view at its peaks.
\item \textbf{Supply--demand coupling and escalation.} Risk arises when supply falls below demand: the same 2\,s glance costs nothing on an empty highway and matters at a crosswalk. The fast loop escalates from a visual cue to a chime to haptics as the gap persists, with a tolerance that shortens as demand rises.
\item \textbf{Environment and sensing degradation.} Lighting, rain and glare raise scene demand and degrade the measurements themselves, gaze estimation and exterior perception alike (Figures~\ref{fig:weather} and~\ref{fig:perception}).
\item \textbf{Adaptation over trips.} Each driver--vehicle pair has a baseline learned over its first trips; the slow loop weighs events against it and recalibrates thresholds over weeks, replaying rare high-risk scenes in CARLA before an over-the-air update (Figure~\ref{fig:highrisk}).
\end{enumerate}

\begin{figure}[H]\centering\includegraphics[width=0.9\linewidth]{driver_view.png}
\caption{The driver's view from the eye point (CARLA, Tesla Model~3 cabin) with the gaze zones a cabin camera classifies. Off-road time is the time the gaze spends outside the green zone; its cost depends on what is in the green zone at that moment.}\label{fig:driver}\end{figure}

\begin{figure}[H]\centering
\begin{subfigure}{0.46\linewidth}\includegraphics[width=\linewidth]{scene_state.png}\caption{Environment state at one moment: ego, nearest pedestrian, lead vehicle, agents within 50\,m, junction and signal, and the demand index they give.}\label{fig:scene}\end{subfigure}\hfill
\begin{subfigure}{0.52\linewidth}\includegraphics[width=\linewidth]{perception.png}\caption{What exterior perception returns: camera image, semantic classes, range per pixel. On the road these are the sources of the scene-demand quantities; in CARLA they are exact.}\label{fig:perception}\end{subfigure}
\caption{Environment state as the twin needs it.}\end{figure}

\begin{figure}[H]\centering\includegraphics[width=0.95\linewidth]{timeseries.png}
\caption{Two minutes of one vehicle in CARLA traffic: vehicle state (speed, accelerations with the 0.3\,g hard-braking line), lead-vehicle headway and TTC, nearest pedestrian ahead and agents within 50\,m, and the illustrative scene-demand index. Shaded bands are junctions; dashed lines mark the moments in Figure~\ref{fig:moments}.}\label{fig:timeseries}\end{figure}

\begin{figure}[H]\centering\includegraphics[width=\linewidth]{demand_moments.png}
\caption{The driver's view at the four most demanding moments of the run, with the quantities that made them demanding. The same glance away from the road would be scored differently in each.}\label{fig:moments}\end{figure}

\begin{figure}[H]\centering\includegraphics[width=\linewidth]{weather.png}
\caption{Same street under clear noon, sunset, hard rain and night: demand rises and sensing degrades while the ground truth is unchanged.}\label{fig:weather}\end{figure}

\begin{figure}[H]\centering\includegraphics[width=0.58\linewidth]{rift_plans_frame.png}
\caption{A generated high-risk scene for calibration: a learned traffic policy drives the marked cars toward the ego (planned paths shown), producing the rare near-miss events that are hard to collect on the road.}\label{fig:highrisk}\end{figure}

\section{Key performance variables}
\begin{table}[H]\footnotesize\centering
\begin{tabularx}{\linewidth}{@{}l X l l@{}}\toprule
\textbf{Variable} & \textbf{Meaning} & \textbf{Unit} & \textbf{Typical / threshold}\\\midrule
\multicolumn{4}{@{}l}{\emph{Driver state}}\\
Gaze zone & which zone the eyes are on & class & 7 zones (Fig.~\ref{fig:driver})\\
Off-road glance duration & continuous time eyes off the road zone & s & risk doubles above 2\\
Eyes-off-road ratio & off-road time over a 10\,s window & \% & $<$30 normal\\
PERCLOS & fraction of time eyelids $>$80\% closed, 1\,min & \% & $>$15 drowsy\\
Head yaw / pitch & head orientation & deg & $\pm$60 / $\pm$30\\
Hands-on torque & steering-column torque from the hands & N\,m & $>$0.3 hands on\\\midrule
\multicolumn{4}{@{}l}{\emph{Vehicle state}}\\
Speed & vehicle speed & m/s (km/h) & 0 to 36 (130)\\
Longitudinal accel. & hard braking threshold 0.3\,g & m/s$^2$ & $<-2.9$ counts\\
Lateral accel. & aggressive turning threshold 0.4\,g & m/s$^2$ & $>$3.9 counts\\
Steering angle / rate & driver or Autosteer input & deg, deg/s & \\
Automation mode & manual, TACC, Autosteer, FSD & class & \\
Takeover reaction time & alert to hands or eyes back & s & $<$1.5 expected\\\midrule
\multicolumn{4}{@{}l}{\emph{Environment state}}\\
Following gap / headway & distance and time to the lead vehicle & m, s & $<$1\,s unsafe following\\
Time-to-collision & gap over closing speed & s & $<$1.5 critical\\
Nearest pedestrian & distance to the closest pedestrian ahead & m & $<$20 demanding\\
Agents within 50\,m & vehicles and pedestrians & count & 0 to 20\\
Junction / signal & in junction, light state & bool, \{R,Y,G\} & \\
Speed limit, lane width & from the map & m/s, m & \\
Weather, illuminance & rain, fog, sun angle, light & class, lux & \\\midrule
\multicolumn{4}{@{}l}{\emph{Twin outputs}}\\
Scene demand index & attention the moment requires & 0 to 1 & Fig.~\ref{fig:timeseries}\\
Attention risk score & demand minus supply, baseline-weighted & 0 to 1 & alert above threshold\\
Alert lead time & warning before a near-miss & s & target $>$ fixed-timer baseline\\
30-day score & context-weighted, explainable & 0 to 100 & mirrors the Safety Score\\
\bottomrule\end{tabularx}
\caption{Key performance variables of the driver--vehicle--environment system.}\label{tab:kpv}\end{table}

\section{Temporal and spatial scales}
\begin{table}[H]\footnotesize\centering
\begin{tabularx}{\linewidth}{@{}p{3.4cm} p{5.0cm} X@{}}\toprule
\textbf{Scale} & \textbf{Value} & \textbf{Why it matters}\\\midrule
Cabin camera / gaze & 10 to 30\,Hz & a 2\,s glance threshold needs about 0.1\,s resolution\\
CAN bus & 10 to 100\,Hz & captures a 0.3\,g braking event (0.2 to 0.5\,s)\\
Exterior perception & 10 to 20\,Hz; CARLA tick 20\,Hz & TTC changes by 0.1\,s per tick at closing speed\\
Glance / blink & 0.1 to 3\,s & the unit of attention supply\\
Alert escalation & 1 to 10\,s, context-dependent & fast-loop window\\
Trip & minutes to hours & drowsiness builds; trip summary\\
Score window & 30 days rolling & slow loop; mirrors Tesla's window\\
Fleet recalibration & months & weights and thresholds revised\\\midrule
Cabin & about 1\,m camera to face & gaze-estimation geometry\\
Vehicle & 4.7\,$\times$\,1.85\,m & own footprint for gaps\\
Lane & 3.5\,m & lane-position variation\\
Perception range & about 100\,m (4\,s at 90\,km/h) & how early demand is visible\\
Pedestrian relevance & under 20 to 30\,m & demand rises sharply\\
\bottomrule\end{tabularx}
\caption{Temporal and spatial scales.}\label{tab:scales}\end{table}

\section{Hand-over to Parts 2.2 and 2.3}
\begin{multicols}{2}\footnotesize
\begin{itemize}
\item \textbf{Driver}: cabin IR camera 10 to 30\,Hz (gaze zone, head pose, PERCLOS from DMD / Drive\&Act-trained models); steering torque 50\,Hz; seat and belt switches.
\item \textbf{Vehicle}: CAN at 10 to 100\,Hz: speed, accelerations, steering, brake, throttle, automation mode, alerts issued and driver responses.
\item \textbf{Environment}: exterior cameras 20 to 36\,Hz, perception at 10 to 20\,Hz (agents, lanes, signals, TTC); map (speed limit, junctions); weather API.
\item \textbf{CARLA}: synchronized ground truth of every agent at 20\,Hz, with weather and generated high-risk scenes; \texttt{attention\_sheet.py} logs every scene-demand quantity per tick and produced Figures~\ref{fig:driver} to~\ref{fig:moments}.
\end{itemize}
\end{multicols}
\vspace{-8pt}
{\small
\noindent\textbf{Part 2.2.} Use Figure~\ref{fig:layout} as the sensor layout. Justify rates against Table~\ref{tab:scales}. Data path: cabin video is processed on the vehicle into features; features, CAN and perception are time-aligned on the vehicle clock and only derived features leave the car. Uncertainty: gaze error with sunglasses, low light and head turns; perception error in rain and at night (Figure~\ref{fig:weather}), quantifiable in CARLA against exact state; CAN noise is small; missing data when the face is not visible. Pre-processing: gaze-zone classification, glance segmentation, PERCLOS over 1\,min, CAN resampling to 10\,Hz, TTC and headway from perception, demand index from Table~\ref{tab:kpv}.

\noindent\textbf{Part 2.3.} Rows come from Table~\ref{tab:kpv} grouped as driver, vehicle, environment and twin outputs; the stream list gives sensor, unit and rate. Tag each row measured, derived or (in CARLA) ground truth. Owners: driver-monitoring team (cabin), vehicle software (CAN, alerts), perception team (environment), OEM safety engineering (score, calibration). A data-flow diagram can follow Figure~\ref{fig:coupling}.\par}
\end{document}
